\documentclass[10pt,a4paper]{amsart}
\usepackage[margin=19mm]{geometry}
\usepackage{amsmath,amssymb,mathtools}
\usepackage[scr=rsfs,cal=euler]{mathalfa}
\usepackage{xfrac}
\usepackage[unicode,hidelinks,bookmarks=false]{hyperref}
\newcommand{\ie}{i.e.\ }
\newcommand{\cf}{cf.\ }
\newcommand{\resp}{respectively\ }
\newcommand{\Subsection}{\subsection}
\providecommand{\nobreakdash}{\nobreak\hbox{-}}
\setlength{\emergencystretch}{2em}
\multlinegap0pt
\usepackage{amssymb}
\usepackage[shortlabels]{enumitem}
\usepackage{mathtools}
\usepackage{xcolor}
\usepackage{bbm}
\newtheorem{lemma}{Lemma}
\newtheorem{proposition}{Proposition}
\def \R {{\mathbb{R} }}
\def \d {{\rm{d}}}
\def \pt{\partial_{t}}
\title{Generically sharp decay and blowing up at infinity for a weak null wave system}
\author{Shijie Dong, Siyuan Ma, Yue Ma, Xu Yuan}
\date{Source: arXiv:2602.22749v1 (CC BY 4.0)}
\begin{document}
\maketitle

\noindent\emph{Source and licence.} Generically sharp decay and blowing up at infinity for a weak null wave system. Version arXiv:2602.22749v1 (CC BY 4.0), distributed by arXiv under the Creative Commons Attribution 4.0 International licence, http://creativecommons.org/licenses/by/4.0/, as recorded in the arXiv licence metadata for this exact version and shown on the abstract page https://arxiv.org/abs/2602.22749v1. Full licence notice: \texttt{licenses/arxiv-2602.22749-CC-BY-4.0.txt}.\par\smallskip

\noindent\textit{Editorial scope.} Counted unit: the article's proposition prop:existence (source0.tex lines 945--961). The article's complete original proof of it is delivered with it (source0.tex lines 963--1177, one \texttt{proof} environment of 215 lines). No mathematics is written, repaired or re-derived: the statement and the proof are the article's own lines, in the article's own notation, and no retained line is substituted or re-flowed.  Retained uncounted as the source-local context the proof needs: lines 135--142; lines 211--219; lines 225--231; lines 580--584; lines 623--626; lines 669--676; lines 687--698; lines 701--715; lines 770--828.  Nothing was omitted from any retained span, so the package carries no omission record.  No retained line contains a citation, so the wrapper carries no bibliography and no bibliography entry.  No retained line contains a figure environment, an image-inclusion command or drawing code, so no image is delivered and the package declares no asset dependency.  Editorial formatting only, in addition: an AMS \texttt{amsart} preamble that restates the article's own declarations and macros used by the retained lines, namely the environments \texttt{lemma}, \texttt{proposition} on separate counters, together with the standard packages those lines need.  The retained environments print in this document as Lemma 1, Proposition 1 in order of appearance, whereas the article prints the same items with the numbering of its own full document.  The complete native source of the article as distributed in the arXiv e-print for this exact version is bundled unchanged as \texttt{sources/195347/source0.tex}.
 \begin{equation}\label{equ:main}
    \left\{
    \begin{aligned}
        -\Box \phi&=\left(\pt \psi\right)^{2},\quad \ (t, x)\in \{ (t,x):t^2 -r^2 \geq 1 \},
        \\
        -\Box \psi&=Q_{0}(\phi,\phi),\ \  (t, x)\in \{ (t,x):t^2 -r^2 \geq 1 \},
    \end{aligned}\right.
 \end{equation}

\begin{equation}\label{equ:main-new}
    \left\{
    \begin{aligned}
        -\Box \phi&=\left(\pt \psi\right)^{2},
        \\
        -\Box \bar{\psi} &= \phi (\partial_t \psi)^2.
        \\
    \end{aligned}\right.
 \end{equation}

\begin{equation}\label{est:smallness1}
\begin{aligned}
    \|\langle x\rangle^{{N}+2} \partial_x^{\leq N+1} (\phi_0,  \psi_0)\| 
    +  \|\langle x\rangle^{{N}+2} \partial_x^{\leq N} (\phi_1, \psi_1)\|   
    <\epsilon,
    \end{aligned}
\end{equation}

\begin{equation}\label{est:t-rpoint}
   u \left|\partial f\right|+u\left|Uf\right|
   +
    v|Vf|\lesssim |\Gamma f|.
\end{equation}

 \begin{equation}\label{est:L2energy}
     \|(s/t)\partial f\|_{L^{2}(\mathcal{H}_{s})}+\|Vf\|_{L^{2}(\mathcal{H}_{s})}+
      \|(1/t)L f\|_{L^{2}(\mathcal{H}_{s})}\lesssim \mathcal{E}(s,f)^{\frac{1}{2}}.
 \end{equation}

 \begin{lemma}
     [Klainerman-Sobolev inequality]\label{le:Klain}
     Let $f=f(t,x)$ be a {sufficiently }regular function defined on $\mathcal{H}_{s}$ with $s\in  [1,+\infty)$. Then we have 
     \begin{equation*}
         \sup_{\mathcal{H}_{s}}\big|t^{\frac{3}{2}}f(t,x)\big|
         \lesssim \sum_{|K|\le 2}\left\|L^{K}f\right\|_{L^{2}(\mathcal{H}_{s})}.
     \end{equation*}
 \end{lemma}

 \begin{lemma}[Estimate for the null form {$Q_0$}]\label{le:null}
     Let $f=f(t,x)$ and $g=g(t,x)$ be two sufficiently regular functions defined on $\mathcal{H}_{s}$ with $s\in [1,+\infty)$. Then we have 
     \begin{equation*}
         \left|Q_{0}(f,g)\right|\lesssim 
         (s/t)^{2}
         |\partial_{t}f||\partial_{t}g|
         +t^{-1}
         \left(\left|Lf\right|\left|\partial g\right|
         +\left|\partial f\right|\left|L g\right|
         \right).
     \end{equation*}
 \end{lemma}

 \begin{lemma}\label{le:energy}
 Let $f=f(t,x)$ be the solution to the Cauchy problem
 \begin{equation*}
     -\Box f=F\quad \mbox{with}\quad (f,\pt f)_{|\mathcal{H}_{1}}=(f_{0},f_{1}).
     \end{equation*}
 Then for any $s\geq 1$, we have 
 \begin{equation*}
 \begin{aligned}
     \mathcal{E}(s,f)^{\frac{1}{2}} &\lesssim
     \mathcal{E}(1,f)^{\frac{1}{2}}+\int_{1}^{s}\|F\|_{L^{2}(\mathcal{H}_{\tau})}\mathrm{d} \tau,\\
     \mathcal{E}_{\rm{con}}(s,f)^{\frac{1}{2}} &\lesssim
     \mathcal{E}_{\rm{con}}(1,f)^{\frac{1}{2}}+\int_{1}^{s}\tau\|F\|_{L^{2}(\mathcal{H}_{\tau})}\mathrm{d} \tau.
     \end{aligned}
 \end{equation*}
     \end{lemma}

 \begin{lemma}\label{le:Linfty}
 The following estimates for the 3D linear wave equations hold.
 \begin{enumerate}
     \item \emph{({Contribution from the initial data}).}
     Let $f=f(t,x)$ be the solution to the Cauchy problem with smooth data
  \begin{equation}\label{equ:waveinitial}
     -\Box f=0\quad \mbox{with}\ \ (f,\pt f)_{|\mathcal{H}_{1}}=(f_{0},f_{1}).
 \end{equation}
Assume that the initial data $(f_{0},f_{1})$ satisfy
\begin{equation}\label{est:smallinitial}
    |f_{0}|+\langle r\rangle\left(|\nabla f_{0}|+|f_{1}|\right)\lesssim {}\langle  r\rangle^{-\frac{7}{2}} \ \ \mbox{and}\ \ 
    \sum_{|I|\le 2}\left\|L^{I}\partial^{\leq 1} f_0\right\|\lesssim 1.
\end{equation}
 Then we have 
 \begin{equation*}
     |f(t,x)|\lesssim (1+u)^{-\frac{3}{2}}(1+v)^{-1}.
 \end{equation*}

\item \emph{({Contribution from the source term}).}
Let $f=f(t,x)$ be the solution to the Cauchy problem
	\begin{equation}\label{equ:wavesource}
	-\Box f=F\quad \mbox{with}\quad (f,\pt f)_{|\mathcal{H}_{1}}=(0,0).
    \end{equation}
	Assume that $F$ is supported in $\mathcal{H}_{[1,+\infty)}$ and satisfies
	\begin{equation*}
	\qquad |F|\le C_{F} (1+u)^{-\mu}(1+v)^{-\nu} \ln^{\gamma} (1+v),\ \  (\mu, \nu, \gamma)\in \mathbb{R}_{+}\times \mathbb{R}_+\times \left\{0,1\right\}.
	\end{equation*}
	Define $F_{\mu}(s)=1,\ln (2+s), (1+s)^{1-\mu}/(1-\mu)$ according to $\mu >1,=1,<1$, respectively, then the following pointwise estimates are true.
	\begin{enumerate}

     \item [\emph{(a)}]  Let $(\mu,\nu,\gamma)\in (1,+\infty)\times \left\{3\right\}\times \left\{1\right\}$. Then we have 
		\begin{equation*}
		|f(t,x)|\lesssim C_{F}(1+u)^{-1}(1+v)^{-1}\ln (2+u).
		\end{equation*}
		
		\item [\emph{(b)}]
        Let $(\mu,\nu,\gamma)\in \R_{+}\times (0,2)\times\left\{0\right\}$.
        Then we have 
		\begin{equation*}
		|f(t,x)|\lesssim C_{F}(1+v)^{-\nu+1}F_{\mu}(u).
		\end{equation*}
		
		\item [\emph{(c)}] 
        Let $(\mu,\nu,\gamma)\in \R_{+}\times \left\{2\right\}\times \left\{0\right\}$.
        Then we have 
		\begin{equation*}
		|f(t,x)|\lesssim C_{F}(1+v)^{-1}\ln (1+v)F_{\mu}(u).
		\end{equation*}
		
		\item [\emph{(d)}]
        Let $(\mu,\nu,\gamma)\in \mathbb{R}_{+}\times (2,+\infty)\times \left\{0\right\}$.
        Then we have 
		\begin{equation*}
		|f(t,x)|\lesssim C_{F}(1+u)^{-(\nu-2)}(1+v)^{-1}F_{\mu}(u).
		\end{equation*}
\end{enumerate}
 
 \end{enumerate}
 \end{lemma}

\begin{proposition}\label{prop:existence}
 Let $N\in \mathbb{N}^{+}$ with $N\ge 6$.
 Let $(I_{1},I_{2})\in \mathbb{N}^{11}$ with $|I_{1}|\le N-3$ and $|I_{2}|\le N-4$.
 There exists an $\epsilon_{0}>0$ such that for all $\epsilon\in (0,\epsilon_{0})$ and all initial data $(\phi_{0},\phi_1,\psi_{0},\psi_1)$ on $\mathcal{H}_{1}$ satisfying the smallness condition~\eqref{est:smallness1}, 
 the Cauchy problem~\eqref{equ:main} admits a global-in-time solution $(\phi,\psi)$ which enjoys the following almost sharp decay estimates 
 \begin{equation*}
 \begin{aligned}
\left|\Gamma^{I_{1}}\phi\right|&\lesssim_{N} \epsilon (1+v)^{-1}\ln(1+v), \ \
\left|\partial\Gamma^{I_{2}}\phi\right|\lesssim_{N} \epsilon (1+v)^{-1}(2+u)^{-1}\ln (1+v),\\
\left|V\Gamma^{I_{2}}\phi\right|&\lesssim_{N}\epsilon (1+v)^{-2}\ln (1+v),\ \
\left|U\Gamma^{I_{2}}\phi\right|\lesssim_{N} \epsilon (1+v)^{-1}(2+u)^{-1}\ln (1+v),\\
\left|\Gamma^{I_{1}}\psi\right|&\lesssim_{N} \epsilon  (1+v)^{-1}(2+u)^{-1}\left(\ln (2+u) + ((2+u)/(1+v))\ln^{2} (1+v)\right),\\
\left|\partial\Gamma^{I_{2}}\psi\right|&\lesssim_{N} \epsilon  (1+v)^{-1}(2+u)^{-2}\left(\ln (2+u) + ((2+u)/(1+v))\ln^{2} (1+v)\right),\\
\left|V\Gamma^{I_{2}}\psi\right|&\lesssim_{N} \epsilon  (1+v)^{-2}(2+u)^{-1}\left(\ln (2+u) + ((2+u)/(1+v))\ln^{2} (1+v)\right).
     \end{aligned}
 \end{equation*}
\end{proposition}

\begin{proof}
 The proof of Proposition~\ref{prop:existence} relies on a bootstrap argument for a high-order energy of the solution $(\phi,\psi)$. From the definition of $\mathcal{E}(s,f)$, the smallness condition~\eqref{est:smallness1} implies the smallness condition for a high-order energy over the initial hypersurface $\mathcal{H}_{1}$:
\begin{equation}\label{est:smallness}
     \sum_{|I|\le N}\mathcal{E}(1,\Gamma^{I}\phi)^{\frac{1}{2}}+\sum_{|I|\le N}\mathcal{E}(1,\Gamma^{I}\psi)^{\frac{1}{2}}
     \lesssim \epsilon.
\end{equation}
In addition, we consider the following bootstrap assumption for the high-order energy over $\mathcal{H}_{s}$: for $C\gg 1$ and $0<\epsilon\ll C^{-1}$ to be chosen later, 
 \begin{equation}\label{est:Boot}
     (\ln(1+s))^{-1}\sum_{|I|\le N}\mathcal{E}(s,\Gamma^{I}\phi)^{\frac{1}{2}}+\sum_{|I|\le N}\mathcal{E}(s,\Gamma^{I}\psi)^{\frac{1}{2}}
     \leq C\epsilon.
 \end{equation}
 For all initial data $(\phi_{0},\phi_1,\psi_{0},\psi_1)$ satisfying~\eqref{est:smallness1}, we set 
\begin{equation*}
S_{*}=S_{*}(\phi_{0},\phi_1,\psi_{0},\psi_1)=\sup \left\{s\in [1,+\infty):(\phi,\psi) \ \mbox{satisfies}~\eqref{est:Boot}\ \mbox{on}\ [1,S_{*})\right\}>1.
\end{equation*}

We are in a position to complete the proof of 
Proposition~\ref{prop:existence}. Note that, in the following discussions, the implicit constants in $\lesssim$ can depend on the constant $N\in \mathbb{N}^{+}$.

    First of all, from~\eqref{est:L2energy}, Lemma~\ref{le:Klain}, the bootstrap assumption \eqref{est:Boot}, and the definition of the energy $\mathcal{E}(s,f)$,  we deduce, for all $s\in [1,S_{*})$, 
\begin{equation}\label{est:Bootpoint}
    \begin{aligned}
       \sum_{|I|\le N-2}
       \left(\left|\partial \Gamma^{I}\psi\right|+s^{-1}\left|L\Gamma^{I}\psi\right|\right)
      & \lesssim C\epsilon t^{-\frac{1}{2}}s^{-1},\\
      (\ln(1+s))^{-1}
      \sum_{|I|\le N-2}
       \left(\left|\partial \Gamma^{I}\phi\right|+s^{-1}\left|L\Gamma^{I}\phi\right|\right)
      & \lesssim C\epsilon t^{-\frac{1}{2}}s^{-1}.
       \end{aligned}
    \end{equation}
    On the other hand, from~\eqref{equ:main}, for any $I\in \mathbb{N}^{11}$ with $|I|\le N$, there exist constants $\left\{a_{J}^{I}\right\}_{J=0}^{I}$  (independent of $(\phi,\psi)$) such that
\begin{equation*}
-\Box{\Gamma}^{I}\phi=a_{J}^{I}{\Gamma}^{J}\left(\partial_{t}\psi\right)^{2} \ \ \mbox{and} \ \
-\Box {\Gamma}^{I}\psi={ a}_{J}^{I}{\Gamma}^{J}Q_{0}(\phi,\phi).
\end{equation*}
Here, we have used the following commutator relations: 
\begin{equation*}
[-\Box, \partial]=[-\Box, L]=[-\Box, \Omega]=[-\Box,S]+2\Box=0.
\end{equation*}
The above identities will be frequently used in this proof. 

\smallskip
{\bf Step 1.} Global existence. We start with the proof of $S_{*}=+\infty$. 
For any initial data $(\phi_{0},\phi_1,{\psi}_{0},\psi_1)$ satisfying the smallness condition~\eqref{est:smallness1}, we consider the corresponding solution $(\phi,\psi)$ of~\eqref{equ:main}. In what follows, we prove the global existence of $(\phi,\psi)$ by improving the estimates of $(\phi,\psi)$ in the bootstrap assumption~\eqref{est:Boot}.

First, from the smallness condition~\eqref{est:smallness1} and Lemma~\ref{le:Klain}, we directly have 
\begin{equation}\label{est:initsmall}
\begin{aligned}
    \sum_{|I|\le N-3}\left(|\Gamma \Gamma^{I}\phi|+|\Gamma\Gamma^{I}\psi|\right)_{|\mathcal{H}_{1}}\lesssim \epsilon\langle r\rangle^{-\frac{5}{2}},\\
    \sum_{|I|\le N-3}
    \left(|\partial\Gamma \Gamma^{I}\phi|+|\partial\Gamma\Gamma^{I}\psi|\right)_{|\mathcal{H}_{1}}
    \lesssim \epsilon\langle r\rangle^{-\frac{7}{2}}.
    \end{aligned}
\end{equation}
For the case of $u\in (0,1)$, from~\eqref{est:Bootpoint} and the definition of $s$, 
\begin{equation*}
   \sum_{|I|\le N-2}
       \left|U\Gamma^{I}\psi\right|\lesssim  \sum_{|I|\le N-2}
       \left|\partial \Gamma^{I}\psi\right|\lesssim C\epsilon  v^{-1}u^{-\frac{1}{2}}.
\end{equation*}
Integrating the above estimate from $(v^{-1},v,\omega)$ which belongs to the initial hypersurface $\mathcal{H}_{1}$, and then using~\eqref{est:initsmall}, we obtain
\begin{equation}\label{est:localu}
    \sum_{|I|\le N-3}\left|\partial \Gamma^{I}\psi\right|\lesssim \epsilon\langle v-v^{-1}\rangle^{-\frac{5}{2}}+
    C\epsilon v^{-1}\int_{v^{-1}}^{u}\tau^{-\frac{1}{2}}\d \tau \lesssim 
    C\epsilon v^{-1}u^{\frac{1}{2}},
\end{equation}
which directly implies 
\begin{equation*}
    \sum_{|I|\le N-3}\left|t\partial \Gamma^{I}\psi\right|\lesssim C\epsilon tv^{-1}u^{\frac{1}{2}}\lesssim  C\epsilon \Longrightarrow 
    \sum_{|I|\le N-3}\|(t/s)\partial\Gamma^{I}\psi\|_{L^{\infty}(\mathcal{H}_{s})}\lesssim C\epsilon s^{-1}.
\end{equation*}
For the case of $u\in (1,+\infty)$, using again~\eqref{est:Bootpoint} and the definition of $s$, we infer
\begin{equation*}
    \sum_{|I|\le N-3}\left|t\partial \Gamma^{I}\psi\right|\lesssim C\epsilon t^{\frac{1}{2}}s^{-1}\lesssim  C\epsilon \Longrightarrow 
    \sum_{|I|\le N-3}\|(t/s)\partial\Gamma^{I}\psi\|_{L^{\infty}(\mathcal{H}_{s})}\lesssim C\epsilon s^{-1}.
\end{equation*}
Combining the above estimates with~\eqref{est:L2energy},~\eqref{est:Boot} and $N\ge 6$, we have 
\begin{equation*}
\begin{aligned}
    \sum_{|I|\le N}\|\Gamma^{I}(\pt \psi)^{2}\|_{L^{2}(\mathcal{H}_{s})}
    &\lesssim \sum_{\substack{|I_{1}|\le N-3\\ |I_{2}|\le N}}
    \|(t/s)\partial\Gamma^{I_{1}}\psi\|_{L^{\infty}(\mathcal{H}_{s})}\|(s/t)\partial\Gamma^{I_{2}}\psi\|_{L^{2}(\mathcal{H}_{s})}\\
    &\lesssim \sum_{\substack{|I_{1}|\le N-2\\ |I_{2}|\le N}}
    \|(t/s)\partial\Gamma^{I_{1}}\psi\|_{L^{\infty}(\mathcal{H}_{s})}\mathcal{E}(s,\Gamma^{I_{2}}\psi)\lesssim C^{2}\epsilon^{2}s^{-1}.
    \end{aligned}
\end{equation*}
Integrating the above estimate over $[1,s]$ and then using Lemma~\ref{le:energy}, we obtain
\begin{equation*}
    \sum_{|I|\le N}\mathcal{E}(s,\Gamma^{I}\phi)^{\frac{1}{2}}
    \lesssim  \epsilon
    +\sum_{|I|\le N}\int_{1}^{s}\|\Gamma^{I}(\pt \psi)^{2}\|_{L^{2}(\mathcal{H}_{\tau})}\d \tau
    \lesssim
    \epsilon+C^{2}\epsilon^{2}\ln (1+s).
\end{equation*}
This estimate strictly improves the energy estimates of $\phi$ in the bootstrap assumption~\eqref{est:Boot} for $C$ large enough and $\epsilon$ small enough.

\smallskip
Second, using again $N\ge 6$, we have, for any $s\in [1,S_{*})$,  
\begin{equation*}
     \sum_{|I|\le N}\|\Gamma^{I}Q_{0}(\phi,\phi)\|_{L^{2}(\mathcal{H}_{s})}
    \lesssim \sum_{\substack{|I_{1}|\le N-3\\ |I_{2}|\le N}}
    \|Q_{0}(\Gamma^{I_{1}}\phi,\Gamma^{I_{2}}\phi)\|_{L^{2}(\mathcal{H}_{s})}.
\end{equation*}
Based on the above estimate and Lemma~\ref{le:null}, we obtain 
\begin{equation*}
\begin{aligned}
     \sum_{|I|\le N}\|\Gamma^{I}Q_{0}(\phi,\phi)\|_{L^{2}(\mathcal{H}_{s})}
    &\lesssim \sum_{\substack{|I_{1}|\le N-3\\ |I_{2}|\le N}}
    \|\partial \Gamma^{I_{1}}\phi\|_{L^{\infty}(\mathcal{H}_{s})}
    \|t^{-1}L\Gamma^{I_{2}}\phi\|_{L^{2}(\mathcal{H}_{s})}\\
    &+\sum_{\substack{|I_{1}|\le N-3\\ |I_{2}|\le N}}
    \|s^{-1}L\Gamma^{I_{1}}\phi\|_{L^{\infty}(\mathcal{H}_{s})}
    \|(s/t)\partial \Gamma^{I_{2}}\phi\|_{L^{2}(\mathcal{H}_{s})}\\
    &+\sum_{\substack{|I_{1}|\le N-3\\ |I_{2}|\le N}}
    \|(s/t)\partial_{t}\Gamma^{I_{1}}\phi\|_{L^{\infty}(\mathcal{H}_{s})}\|(s/t)\partial_{t}\Gamma^{I_{2}}\phi\|_{L^{2}(\mathcal{H}_{s})},
\end{aligned}
\end{equation*}
which, combined with the estimates ~\eqref{est:L2energy},~\eqref{est:Boot} and~\eqref{est:Bootpoint}, implies
\begin{equation*}
    \sum_{|I|\le N}\|\Gamma^{I}Q_{0}(\phi,\phi)\|_{L^{2}(\mathcal{H}_{s})}\lesssim C^{2}\epsilon^{2}s^{-\frac{3}{2}}
    \ln^{2} (1+s)\lesssim C^{2}\epsilon^{2}s^{-\frac{5}{4}}.
\end{equation*}
Integrating the above estimate over $[1,s]$, and then using Lemma~\ref{le:energy}, we deduce
\begin{equation*}
    \sum_{|I|\le N}\mathcal{E}(s,\Gamma^{I}\psi)^{\frac{1}{2}}
    \lesssim  \epsilon
    +\sum_{|I|\le N}\int_{s_{0}}^{s}\|\Gamma^{I}Q_{0}(\phi,\phi)\|_{L^{2}(\mathcal{H}_{\tau})}\d \tau
    \lesssim
    \epsilon+C^{2}\epsilon^{2}.
\end{equation*}
This estimate strictly improves the energy estimates of $\psi$ in the bootstrap assumption~\eqref{est:Boot} for  $C$ large enough and $\epsilon$ small enough.

\smallskip
We have so far strictly improved the estimates of $(\phi,\psi)$ in the bootstrap assumption~\eqref{est:Boot} and, thus, for all initial data $(\phi_{0},\phi_1, {\psi}_{0},\psi_1)$ satisfying~\eqref{est:smallness1}, {$S_{*}(\phi_{0},\psi_{0})=+\infty$ and} the proof of global existence is complete.

\smallskip
\textbf{Step 2.} Almost sharp decay estimates.
In what follows, we tacitly write the decay rates of the solution pair in terms of $(u,v)$,
while by the local results (see, for example, the proof of~\eqref{est:localu}), the same decay estimates remain valid if one replaces $(u,v)$ by $(2+u,1+v)$.

\smallskip
From~\eqref{est:Bootpoint}, for any $J\in \mathbb{N}^{11}$ with $|J|\le N-2$, we directly have 
\begin{equation*}
    \left|\Gamma^{J}(\pt \psi )^{2}\right|
    \lesssim \sum_{\substack{|J_{1}|\le N-2\\ |J_{2}|\le N-2}}
    \left|\partial \Gamma^{J_{1}}\psi\right|
    \left|\partial \Gamma^{J_{2}}\psi\right|\lesssim \epsilon t^{-1}s^{-2}\lesssim \epsilon v^{-2}u^{-1}.
\end{equation*}
Therefore, from Lemma~\ref{le:Linfty},  we obtain,   for any $J\in \mathbb{N}^{+}$ with $|J|\le N-2$,
\begin{equation*}
    \left| \Gamma^{J}\phi\right|\lesssim 
    \epsilon v^{-1}u^{-1}+\epsilon v^{-1}(\ln v) (\ln u)\lesssim \epsilon v^{-1}\ln^{2}v.
\end{equation*}
Here, we have used (i) of Lemma~\ref{le:Linfty} and the smallness condition~\eqref{est:smallness1} for initial data.
Combining the above estimate with~\eqref{est:t-rpoint}, we obtain a rough decay estimate for $\phi$ with an additional $\ln v$ growth.

\smallskip
Next, from the above decay estimates for $\phi$ and Lemma~\ref{le:null}, for any $J\in \mathbb{N}^{11}$ with $|J|\le N-3$, we directly have
\begin{equation*}
\begin{aligned}
     |\Gamma^{J}Q_{0}(\phi,\phi)|
    &\lesssim t^{-1}\sum_{\substack{|J_{1}|\le N-3\\ |J_{2}|\le N-3}}
    |\partial \Gamma^{J_{1}}\phi|
    |L\Gamma^{J_{2}}\phi|\\
    &+t^{-1}\sum_{\substack{|J_{3}|\le N-3\\ |J_{4}|\le N-3}}
    |L\Gamma^{J_{3}}\phi|
    |\partial\Gamma^{J_{4}}\phi|\\
    &+(s/t)^{2}\sum_{\substack{|J_{5}|\le N-3\\ |J_{6}|\le N-3}}
    |\partial\Gamma^{J_{5}}\phi|
    |\partial\Gamma^{J_{6}}\phi|\lesssim \epsilon v^{-3}u^{-1}\ln^{4}v.
\end{aligned}
\end{equation*}
Therefore, from Lemma~\ref{le:Linfty},  we infer, for any $J\in \mathbb{N}^{11}$ with $|J|\le N-3$,
\begin{equation*}
    \left| \Gamma^{J}\psi\right|\lesssim \epsilon v^{-1}u^{-1}+\epsilon v^{-1}u^{-1+\delta_{1}}\lesssim \epsilon v^{-1}u^{-1+\delta_{1}},\quad \mbox{for any}\ \ 0<\delta_{1}\ll 1.
\end{equation*}
Here, we have used again (i) of Lemma~\ref{le:Linfty} and the smallness condition~\eqref{est:smallness1} for initial data.
Combining the above estimate with~\eqref{est:t-rpoint}, we obtain a rough decay estimate for $\psi$ with an additional $u^{\delta_{1}}$ growth.

\smallskip
Then, from the above estimate for $\psi$ and~\eqref{est:t-rpoint}, for any $J\in \mathbb{N}^{+}$ with $|J|\le N-3$, 
\begin{equation*}
    \left| \Gamma^{J}(\pt \psi )^{2}\right|
    \lesssim \sum_{\substack{|J_{1}|\le N-3\\ |J_{2}|\le N-4}}
    \left|\partial \Gamma^{J_{1}}\psi\right|
    \left|\partial \Gamma^{J_{2}}\psi\right|
    \lesssim \epsilon v^{-2}u^{-\frac{5}{2}+\delta_{1}}.
\end{equation*}
It follows from Lemma~\ref{le:Linfty} that 
\begin{equation*}
    \left|\Gamma^{J}\phi\right|
    \lesssim 
    \epsilon v^{-1}u^{-1}+
    \epsilon v^{-1}\ln^{}v
    \lesssim \epsilon v^{-1}\ln v,\quad 
    \mbox{for any} \ \ |J|\le N-3.
\end{equation*}
Combining the above estimate with~\eqref{est:t-rpoint}, we complete the proof of estimates for $\phi$.

\smallskip
Last, recall from \eqref{equ:main-new} that the scalar $\bar{\psi} = \psi + \frac{1}{2}\phi^2$ satisfies $-\Box \bar{\psi}=\phi(\pt \psi)^{2}$.
Using the above estimates for $\phi$ and $\psi$, for $J\in \mathbb{N}^{+}$ with $|J|\le N-3$, we have 
\begin{equation*}
   \left|  \Gamma^{J}\left(\phi(\pt \psi)^{2}\right)\right|\lesssim
   \epsilon^3 u^{-\frac{5}{2}+\delta_1}v^{-3}\ln v,\quad \mbox{for any}\ 0<\delta_{1}\ll 1.
\end{equation*}
It follows from Lemma~\ref{le:Linfty} that 
\begin{equation*}
    \big| \Gamma^{J}\bar{\psi}\big|\lesssim \epsilon u^{-1}v^{-1}\ln u\Longrightarrow 
    \left|\Gamma^{J}\psi\right|
\lesssim \epsilon u^{-1}v^{-1}\ln u + \epsilon v^{-2}\ln^{2} v.
\end{equation*}
Combining the above estimates with~\eqref{est:t-rpoint}, we complete the proof.
\end{proof}

\end{document}
